\section{GÜVENSİZ API ORTAMI VE ZAFİYET SENARYOLARI}
\label{bolum6}

API Fortress sistemi içerisinde geliştirilen Vulnerable API ortamı, REST API güvenlik açıklarının uygulamalı olarak öğrenilmesini sağlamak amacıyla tasarlanmıştır. Bu ortamda çeşitli güvenlik açıkları bilinçli olarak sisteme entegre edilmiş; kullanıcıların saldırı senaryolarını kontrollü bir laboratuvar ortamında deneyimleyebilmesi hedeflenmiştir. Geliştirilen zafiyet senaryoları, OWASP API Security Top 10 kapsamında değerlendirilen modern API güvenlik problemlerini temel almaktadır.

Vulnerable API ortamının temel amacı, kullanıcıların yalnızca teorik güvenlik kavramlarını öğrenmesi değil; aynı zamanda bu güvenlik açıklarının gerçek sistemlerde nasıl ortaya çıktığını, saldırganların sistem davranışını nasıl manipüle ettiğini ve eksik güvenlik kontrollerinin ne tür sonuçlara yol açabileceğini uygulamalı olarak gözlemleyebilmesidir.

Bu ortamda bulunan endpoint’ler gerçek dünya REST API mimarisine benzer şekilde tasarlanmıştır. Kullanıcı kayıt işlemleri, kimlik doğrulama süreçleri, veri erişim mekanizmaları ve CRUD operasyonları normal bir backend sistemi mantığıyla çalışmaktadır. Ancak güvenlik kontrolleri bilinçli olarak eksik bırakılmış veya yanlış uygulanmıştır. Böylece kullanıcılar modern API sistemlerinde karşılaşılan zafiyetleri güvenli bir laboratuvar ortamında test edebilmektedir.

\subsection{Mass Assignment}
\label{mass-assignment}

Mass Assignment zafiyeti, istemci tarafından gönderilen verilerin sunucu tarafında yeterli filtreleme yapılmadan doğrudan modele aktarılması sonucunda ortaya çıkan kritik bir güvenlik problemidir. Bu tür zafiyetlerde saldırgan, normalde erişememesi gereken model alanlarını HTTP isteği üzerinden manipüle ederek sistem davranışını değiştirebilir.

Modern REST API sistemlerinde kullanıcıdan gelen JSON verileri çoğu zaman doğrudan veri modellerine dönüştürülmektedir. Eğer geliştirici yalnızca izin verilmesi gereken alanları filtrelemezse, saldırgan ek alanlar göndererek sistem üzerinde yetki yükseltme veya kritik yapılandırma değişiklikleri gerçekleştirebilir.

Bu senaryoda kullanıcı kayıt endpoint’i güvenli olmayan şekilde tasarlanmıştır. Kullanıcıdan gelen JSON verileri herhangi bir filtreleme yapılmadan doğrudan User modeline aktarılmaktadır.

Güvensiz endpoint örneği aşağıdaki gibidir:

\begin{lstlisting}
# Güvensiz kayıt endpoint'i (routes_auth.py)

@auth_bp.route('/register', methods=['POST'])
def register():
    data = request.get_json()

    # Tehlike:
    # Gelen tüm alanlar filtrelenmeden modele aktarılıyor

    new_user = User(**data)

    db.session.add(new_user)
    db.session.commit()

    return jsonify({'message': 'User created'}), 201
\end{lstlisting}

Bu yapı ilk bakışta işlevsel görünse de ciddi güvenlik riski oluşturmaktadır. Çünkü istemciden gelen tüm alanlar doğrudan modele aktarılmaktadır. Sistem yalnızca:

\begin{lstlisting}
{
  "username": "enes",
  "password": "123456"
}
\end{lstlisting}

gibi alanlar beklerken saldırgan isteğe aşağıdaki alanı da ekleyebilir:

\begin{lstlisting}
{
  "username": "attacker",
  "password": "123456",
  "role": "admin"
}
\end{lstlisting}

Bu durumda saldırgan kendi hesabını admin yetkisiyle oluşturabilir. Çünkü backend tarafında hangi alanların kabul edileceği sınırlandırılmamıştır.

\subsubsection{Bu zafiyet özellikle:}

\begin{itemize}[leftmargin=*]
    \item kullanıcı rolleri,
    \item hesap durumu,
    \item erişim seviyeleri,
    \item ödeme bilgileri,
    \item sistem ayarları
\end{itemize}

gibi kritik alanların istemci tarafından değiştirilebilmesine neden olabilir.

Mass Assignment problemi modern API sistemlerinde oldukça yaygın görülen güvenlik açıklarından biridir. Özellikle ORM tabanlı backend sistemlerinde kullanıcıdan gelen verilerin doğrudan modele aktarılması ciddi risk oluşturmaktadır.

\subsubsection{Bu senaryo sayesinde kullanıcılar:}

\begin{itemize}[leftmargin=*]
    \item istemci verisinin neden filtrelenmesi gerektiğini,
    \item güvenli veri doğrulama mantığını,
    \item backend tarafında allowlist yaklaşımının önemini
\end{itemize}

uygulamalı olarak gözlemleyebilmektedir.

\subsection{Broken Authentication}
\label{broken-authentication}

Broken Authentication zafiyeti, kullanıcı kimlik doğrulama süreçlerinde bulunan eksik veya hatalı güvenlik kontrolleri sonucunda ortaya çıkmaktadır. Kimlik doğrulama mekanizmaları modern backend sistemlerinin en kritik bileşenlerinden biridir. Bu mekanizmalardaki küçük bir hata bile saldırganların yetkisiz erişim elde etmesine neden olabilir.

API Fortress içerisindeki Broken Authentication senaryosunda iki farklı güvenlik problemi uygulanmıştır.

\begin{enumerate}[leftmargin=*]
    \item Parola doğrulama süreçlerinin eksik olması

    Bazı durumlarda sistemin parola kontrolünü tam olarak gerçekleştirmemesi, saldırganların geçerli kullanıcı bilgileri olmadan sisteme erişebilmesine neden olabilir.

    \item Gizli HTTP header ile authentication bypass
\end{enumerate}

Bu senaryoda sisteme bilinçli olarak gizli bir arka kapı (backdoor) yerleştirilmiştir.

Güvensiz endpoint aşağıdaki gibidir:

\begin{lstlisting}
# Gizli header ile authentication bypass

@auth_bp.route('/login', methods=['POST'])
def login():

    if request.headers.get('X-Admin-Override') == 'true':
        return jsonify({
            'token': generate_admin_token()
        }), 200

    # Normal akış devam eder...
\end{lstlisting}

Bu senaryoda saldırgan normal kullanıcı bilgilerine ihtiyaç duymadan yalnızca aşağıdaki HTTP header’ını ekleyerek admin token’ı alabilmektedir:

\begin{lstlisting}
X-Admin-Override: true
\end{lstlisting}

\subsubsection{Bu yapı gerçek sistemlerde zaman zaman:}

\begin{itemize}[leftmargin=*]
    \item debug mekanizmaları,
    \item test backdoor’ları,
    \item unutulmuş geliştirme kodları,
    \item yanlış yapılandırılmış admin bypass sistemleri
\end{itemize}

nedeniyle oluşabilmektedir.

\subsubsection{Bu zafiyetin temel amacı:}

\begin{itemize}[leftmargin=*]
    \item authentication mekanizmalarının neden kritik olduğunu,
    \item gizli bypass mekanizmalarının oluşturduğu riski,
    \item backend güvenlik kontrollerinin neden merkezi yönetilmesi gerektiğini
\end{itemize}

uygulamalı olarak göstermektir.

\subsubsection{Broken Authentication problemleri özellikle:}

\begin{itemize}[leftmargin=*]
    \item token manipülasyonu,
    \item session hijacking,
    \item brute force saldırıları,
    \item zayıf parola politikaları,
    \item gizli bypass mekanizmaları
\end{itemize}

gibi güvenlik problemleriyle birlikte modern API sistemlerinde en kritik tehditlerden biri olarak kabul edilmektedir.

\subsection{IDOR / BOLA}
\label{idor-bola}

IDOR (Insecure Direct Object Reference) ve BOLA (Broken Object Level Authorization), modern REST API sistemlerinde en yaygın ve en kritik güvenlik açıklarından biridir.

Bu zafiyetin temel nedeni, sistemin kullanıcıların erişmeye çalıştığı nesneler üzerinde yeterli authorization kontrolü yapmamasıdır.

API Fortress içerisindeki senaryoda kullanıcı detay endpoint’i güvenli olmayan şekilde tasarlanmıştır:

\begin{lstlisting}
# Güvensiz kullanıcı detay endpoint'i

@users_bp.route('/users/<int:user_id>', methods=['GET'])
def get_user(user_id):

    # Tehlike:
    # Kimlik doğrulama yok,
    # herhangi bir ID sorgulanabilir

    user = User.query.get(user_id)

    return jsonify(user.to_dict()), 200
\end{lstlisting}

\subsubsection{Bu endpoint üzerinde:}

\begin{itemize}[leftmargin=*]
    \item JWT doğrulaması bulunmamaktadır,
    \item kullanıcı kontrolü yapılmamaktadır,
    \item erişim yetkisi doğrulanmamaktadır.
\end{itemize}

Bu nedenle saldırgan URL içerisindeki kullanıcı ID’sini değiştirerek farklı kullanıcılara ait verilere erişebilir.

Örneğin:

\begin{verbatim}
GET /users/1
GET /users/2
GET /users/3
\end{verbatim}

şeklinde farklı ID’ler denenerek sistemdeki diğer kullanıcı bilgileri görüntülenebilir.

\subsubsection{Bu saldırı tipi özellikle:}

\begin{itemize}[leftmargin=*]
    \item kullanıcı profilleri,
    \item sipariş sistemleri,
    \item banka işlemleri,
    \item müşteri kayıtları,
    \item sağlık sistemleri
\end{itemize}

gibi hassas veri içeren API sistemlerinde ciddi veri sızıntılarına neden olabilmektedir.

OWASP API Security Top 10 listesinde BOLA, en kritik API güvenlik problemlerinden biri olarak değerlendirilmektedir. Çünkü saldırgan çoğu zaman yalnızca URL parametresini değiştirerek yetkisiz veri erişimi sağlayabilmektedir.

\subsubsection{Bu senaryo sayesinde kullanıcılar:}

\begin{itemize}[leftmargin=*]
    \item object-level authorization mantığını, kullanıcı izolasyonunun önemini, erişim kontrolünün neden endpoint seviyesinde uygulanması gerektiğini uygulamalı olarak analiz edebilmektedir.
\end{itemize}

\subsection{Broken Function Level Authorization}
\label{broken-function-level-authorization}

Broken Function Level Authorization zafiyeti, kullanıcıların erişmemesi gereken kritik sistem fonksiyonlarına ulaşabilmesi durumunda ortaya çıkmaktadır.

\subsubsection{Bu problem genellikle:}

\begin{itemize}[leftmargin=*]
    \item admin endpoint’lerinde,
    \item yönetim panellerinde,
    \item kritik CRUD işlemlerinde,
    \item rol tabanlı erişim kontrollerinde
\end{itemize}

yetersiz authorization uygulanması sonucunda oluşmaktadır.

API Fortress içerisindeki güvensiz kullanıcı silme endpoint’i aşağıdaki gibidir:

\begin{lstlisting}
# Güvensiz kullanıcı silme endpoint'i

@users_bp.route('/users/<int:user_id>', methods=['DELETE'])
def delete_user(user_id):

    # Tehlike:
    # Rol kontrolü yok,
    # herkes herhangi kullanıcıyı silebilir

    user = User.query.get(user_id)

    db.session.delete(user)
    db.session.commit()

    return jsonify({
        'message': 'User deleted'
    }), 200
\end{lstlisting}

\subsubsection{Bu endpoint üzerinde:}

\begin{itemize}[leftmargin=*]
    \item kullanıcı rolü kontrol edilmemektedir,
    \item admin doğrulaması yapılmamaktadır,
    \item authorization middleware kullanılmamaktadır.
\end{itemize}

Bu nedenle normal bir kullanıcı aşağıdaki isteği göndererek başka kullanıcıları silebilir: Bu tür güvenlik açıkları gerçek sistemlerde:

\begin{itemize}[leftmargin=*]
    \item kullanıcı hesaplarının silinmesine,
    \item yönetici işlemlerinin ele geçirilmesine,
    \item veri kaybına,
    \item sistem bütünlüğünün bozulmasına
\end{itemize}

neden olabilir.

Broken Function Level Authorization zafiyetleri özellikle mikro servis mimarilerinde oldukça yaygın görülmektedir. Çünkü geliştiriciler çoğu zaman endpoint’in yalnızca frontend tarafından erişileceğini varsaymakta ve backend tarafında yeterli yetki kontrolü uygulamamaktadır.

\subsubsection{Bu senaryo sayesinde kullanıcılar:}

\begin{itemize}[leftmargin=*]
    \item rol tabanlı erişim kontrolünün önemini,
    \item admin endpoint’lerinin neden korunması gerektiğini,
    \item authorization süreçlerinin neden backend tarafında uygulanması gerektiğini
\end{itemize}

uygulamalı olarak öğrenebilmektedir.

Sonuç olarak Vulnerable API ortamı içerisinde geliştirilen Mass Assignment, Broken Authentication, IDOR/BOLA ve Broken Function Level Authorization senaryoları; modern REST API sistemlerinde yaygın olarak görülen kritik güvenlik problemlerini temsil etmektedir. Bu senaryolar sayesinde kullanıcılar yalnızca teorik güvenlik kavramlarını öğrenmekle kalmamakta; aynı zamanda gerçek dünya saldırı mantığını kontrollü laboratuvar ortamında uygulamalı olarak deneyimleyebilmektedir.


\clearpage
